\documentclass[dvipdfmx,12pt]{jbook}

\usepackage[
  a5paper, % A5判に設定
  twoside, % 見開きレイアウト
  top=15mm, % 上余白
  bottom=15mm, % 下余白
  inner=15mm, % 内余白
  outer=15mm % 外余白
]{geometry}

\usepackage{hyperref}
\usepackage{makeidx}
\usepackage{amsmath}
\usepackage{amssymb}
\usepackage{amsthm}
\usepackage{bm}
\usepackage{pxrubrica}

\makeindex

\theoremstyle{definition}
\newtheorem{dfn}{定義}
\newtheorem*{dfn*}{定義}

\theoremstyle{plain}
\newtheorem{axm}{公理}
\newtheorem*{axm*}{公理}
\newtheorem{thm}{定理}
\newtheorem*{thm*}{定理}
\newtheorem{col}{系}
\newtheorem*{col*}{系}
\newtheorem{lmm}{補題}
\newtheorem*{lmm*}{補題}
\newtheorem{semi}{準公理}
\newtheorem*{semi*}{準公理}

\renewcommand{\proofname}{\textbf{証明}}

\title{平面幾何学の一階述語理論}
\date{\today}
\author{Tsugu}

\begin{document}
  \pagenumbering{roman}
  \maketitle
  \chapter*{前書き}

これは初等幾何学を系統的に述べようとするものではありません。
作者の私的なまとめです。

一階述語論理の幾何学の公理系としては、タルスキーによるものがありますが、公理系がいわば極めて圧縮されているので、論証的展開の初期段階ほど難しいということになる。
加えて、展開初期の頃はむしろ幾何学的直感が推論を邪魔するのである。

一方、こちらの公理系は、確かにタルスキーの公理系よりもだいぶ公理が多いですが、まだしも使いやすいと思います。
（それに、「公理」と銘打たれた部分に関しては単に実数の公理系相当物であり、そこは言ってみれば「常識」で、意味があるのは「公準」と題した部分なのである。）

しかしながら、一階述語論理の理論に宿命的なものとしての使いにくさも、やはりあります。
例えば、直線などが論理的対象ではないことです。
直線の関数をこの体系で考えることはできない。
直線を定める二点の関数を考えることはできる。
また、直線を擬似的に主語のように見せかけて命題を言い表すことはできなくもないが、述語に対するまごうかたなき項として扱うことはできない。

角の概念もまた問題です。
角 \(AOB\) と言ったとき、イメージとしては劣角と優角の二つあるわけです。
ところが最も標準的に考えた場合、角は頂点を共有する二つの半直線であり、しかしこれだけでは大小どちらの角であるのか定まらない。
その暗黙の指定は角 \(AOB\) を考えている人間の頭の中にのみある。
そこで、このような曖昧さを払拭すべく、最も即物的かつシンプルに角を定義すれば、三点のなす三角形の角としての角もしくは三点が一直線上にあるときの平角あるいは大きさゼロの角、となります。
そしてこの場合も、角そのものを直接定義することはできず、角の合同関係を通して間接的に定義されるのです。

このように、明らかに存在しているはずなのに一階述語論理では対象物として扱えないというもどかしさがありますが、その代わり、公理系はゲーデルの不完全性定理の意味で「完全」になります。

  \tableofcontents
  \setcounter{page}{1}
  \pagenumbering{arabic}
  \chapter{平面幾何学の一階述語理論}


\section{無定義用語}

\begin{description}
  \item[定項] \(0,\) \(1\)
  \item[原始関数] \(+,\) \(\cdot\)
  \item[原始述語] \(S,\) \(>,\) \(V\)
\end{description}

\(S(x)\) の直感的な意味は、「\(x\) は\textbf{スカラー}\index{すからー@スカラー}である」である。


\section{公理}

これは実質的に\textbf{実閉体}\index{じつへいたい@実閉体}の公理系である。
実数の公理系でも構わない。
実閉体のものにしてあるのは、一階述語論理において理論の完全性が得られるからである。

\(\exists!x(P(x))\) は \(\exists x(P(x)\land\forall y(P(y)\rightarrow x=y))\) の省略記法であり、「ただ一つの \(x\) に対して \(P(x)\) である」を意味する。
また、以下において \(\forall\) は極力省略してある。

\begin{enumerate}
  \item \(S(0)\land S(1)\)
  \item \(S(x)\land S(y)\leftrightarrow S(x+y)\)
  \item \(S(x)\land S(y)\rightarrow\exists!z(x+z=y)\)
  \item \(S(x)\rightarrow x+0=x\)
  \item \(S(x)\land S(y)\leftrightarrow S(x\cdot y)\)
  \item \(S(x)\land S(y)\rightarrow\exists!z(x\cdot z=y)\)
  \item \(S(x)\rightarrow x\cdot 1=x\)
  \item \(S(x)\land S(y)\land S(z)\rightarrow x+(y+z)=(x+y)+z\)
  \item \(S(x)\land S(y)\rightarrow x+y=y+x\)
  \item \(S(x)\land S(y)\land S(z)\rightarrow x\cdot(y\cdot z)=(x\cdot y)\cdot z\)
  \item \(S(x)\land S(y)\rightarrow x\cdot y=y\cdot x\)
  \item \(S(x)\land S(y)\land S(z)\rightarrow x\cdot(y+z)=x\cdot y+x\cdot z\)
  \item \(S(x)\land S(y)\rightarrow x=y\lor x>y\lor y>x\)
  \item \(S(x)\land S(y)\rightarrow x=y\rightarrow\lnot x>y\lor\lnot y>x\)
  \item \(S(x)\land S(y)\rightarrow x>y\rightarrow x\neq y\lor\lnot y>x\)
  \item \(S(x)\land S(y)\land S(z)\rightarrow(x>y\land y>z\rightarrow x>z)\)
  \item \(S(x)\land S(y)\land S(z)\rightarrow(x>y\rightarrow x+z>y+z)\)
  \item \(S(x)\land S(y)\land S(z)\rightarrow(x>y\land z>0\rightarrow x\cdot z>y\cdot z)\)
  \item\label{260425093801} \(S(y)\rightarrow\{y>0\rightarrow\exists x(S(x)\land x\cdot x=y)\}\)
  \item\label{260425093802} 「変数 \(x\) に関する任意の \(n\) 次多項式」を \(\varphi_n(x)\) \((n=1,2,\ldots)\) と表すとき、\\
        \(\exists x(S(x)\land\varphi_{2n-1}(x)=0).\)
\end{enumerate}

公理に基づいて、負数 \(-x,\) 差 \(x-y,\) 商 \(x/y,\) 自乗 \(x^2,\) 平方根 \(\sqrt{x}\) などがただ一つ存在することが示され、定義可能になる。

公理 \ref{260425093801} は公理ではなく\textbf{公理図式}\index{こうりずしき@公理図式}であり、実際は無数の論理式である。
しかし、アルゴリズム的に生成あるいは機械的に判別可能なものなので、「有限の立場」に反しない。

ちなみに、公理 \ref{260425093801} と \ref{260425093802} をひとまとめにして\(\exists x(S(x)\land\varphi_n(x)=0)\) であるとすると、スカラーの全体が複素閉体になってしまい、大小関係を一般的に定義することができなくなる。


\section{公準}

これは座標ベクトルの公理系である。
\textbf{点}\index{てん@点}は\textbf{ベクトル}\index{べくとる@ベクトル}であり、ベクトルは点である。

\begin{enumerate}
  \item \(\lnot S(z)\leftrightarrow\exists xy(S(x)\land S(y)\land V(x,y,z))\)
  \item\label{260425093804} \(S(x)\land S(y)\leftrightarrow \exists zV(x,y,z)\)
  \item\label{260425155601} \(V(x,y,z)\land V(x,y,w)\rightarrow z=w\)
  \item \(V(x,y,z)\land V(u,v,z)\rightarrow x=u\land y=v\)
  \item\label{260425093803} \(\displaystyle V(x,y,w)\rightarrow z+w=w+z=\left(1+\frac{z}{\sqrt{x^2+y^2}}\right)\cdot w\)
  \item \(V(x,y,w)\land V(z\cdot x,z\cdot y,u)\rightarrow z\cdot w=w\cdot z=u\)
  \item \(V(x,y,z)\land V(u,v,w)\land V(x+u,y+v,t)\rightarrow z+w=t\)
  \item \(V(x,y,z)\land V(u,v,w)\land V(x\cdot u,y\cdot v,t)\rightarrow z\cdot w=t\)
\end{enumerate}

% 定義 1

\begin{dfn}[\textbf{順序対}\index{じゅんじょつい@順序対}]\label{260419033901}
  公準 \ref{260425093804} より、任意のスカラー \(x,y\) に対して \(V(x,y,z)\) であるような \(z\) が存在し、公準 \ref{260425155601} より、\(z\) は \(x,y\) に対してただ一つ定まるので、それを
  \[\langle x,y\rangle\]
  と書く。
\end{dfn}

定義 \ref{260419033901} の記法を用いて公準の本質的な部分を書き換えるとこうなる。

\begin{enumerate}
  \setcounter{enumi}{3}
  \item \(\langle x,y\rangle=\langle u,v\rangle\rightarrow x=u\land y=v\)
  \item\label{260425094701} \(\displaystyle z+\langle x,y\rangle=\langle x,y\rangle+z=\left(1+\frac{z}{\sqrt{x^2+y^2}}\right)\cdot\langle x,y\rangle\)
  \item\label{260425094702} \(z\cdot\langle x,y\rangle=\langle x,y\rangle\cdot z=\langle z\cdot x,z\cdot y\rangle\)
  \item \(\langle x,y\rangle+\langle u,v\rangle=\langle x+u,y+v\rangle\)
  \item\label{260425093806} \(\langle x,y\rangle\cdot\langle u,v\rangle=x\cdot u+y\cdot v\)
\end{enumerate}

結局、公準 \ref{260425094701} 以外は \(\langle x,y\rangle\) がベクトルであることを示している。
\ref{260425094701} は、スカラーとベクトルの和が未定義になることを避けるための、いわばダミーの公準である。
\ref{260425094701} はスカラーの分だけベクトルを延長する。
これは \ref{260425094702} の内容に合わせたものである。

この公理系と諸公準は、\(S(x)\) を「\(x\) は実数である」、\(V(x,y,z)\) を「\(z\) はベクトル \(\langle x,y\rangle\) である」と解釈し、\(+\) を和、\(\cdot\) を積、ただし公準 \ref{260425093806} のそれは内積、と解釈すれば公準 \ref{260425094701} 以外すべて成り立つ。
\ref{260425094701} については、スカラーを加えることをベクトルの延伸と解すれば成り立つ。
ゆえに、集合論が無矛盾ならばこの体系も無矛盾である。

公理と公準より、ベクトルの差 \(\langle x,y\rangle-\langle u,v\rangle\) が定義できる。

% 定義 2

\begin{dfn}[\textbf{ノルム}\index{のるむ@ノルム}]
  \(x\) のノルムを
  \[||x||=\sqrt{x\cdot x}\]
  と定義する。
\end{dfn}

三角形は一直線上にない三点と同一視され、角は三角形のそれに限られる。
角の大小は、角を挟む一方の辺と対辺のそれぞれを重ね合わせた場合における対辺の長短と同一視される。


\section{二次元ユークリッド空間であることの確認}

以下においては、\(S(x)\) を満たすもの、すなわち\textbf{スカラーをギリシャ文字での小文字で表し}、そうでないもの、すなわち\textbf{ベクトルをラテン文字の小文字で表す}。
ベクトルは点と同一視される。

また、演算子 \(\cdot\) はしばしば省略される。

% 定義 3

\begin{dfn}
  任意の二点 \(x,y\) に対して、
  \[y-x\]
  を
  \[\overrightarrow{xy}\]
  と書く。
\end{dfn}

% 定理 1

\begin{thm}
  少なくとも一つの点が存在する。
\end{thm}

\begin{proof}
  例えば \(\langle0,0\rangle.\)
\end{proof}

% 定理 2

\begin{thm}
  任意の点 \(x\) と \(y\) に対して、
  \[\overrightarrow{xy}\]
  がただ一つ存在する。
\end{thm}

\begin{proof}
  \(-\) の定義（未記載）と公理 3 より明らか。
\end{proof}

% 定理 3

\begin{thm}
  任意の点 \(x\) と \(a\) に対して、
  \[\overrightarrow{xy}=a\]
  となる点 \(y\) がただ一つ存在する。
\end{thm}

\begin{proof}
  \(y-x=a\) より \(y=x+a.\)
\end{proof}

% 定理 4

\begin{thm}
  \[\overrightarrow{xy}+\overrightarrow{yz}=\overrightarrow{xz}.\]
\end{thm}

\begin{proof}
  \((y-x)+(z-y)=z-x.\)
\end{proof}

% 定理 5

\begin{thm}
  少なくとも二つの点が存在し、一方は他方の点とスカラーとの積で表せない。
\end{thm}

\begin{proof}
  例えば、\(x=\langle1,0\rangle,\ y=\langle0,1\rangle\) とおくと、\(\lambda\cdot x=\langle\lambda,0\rangle\) であるので常に \(y\neq\lambda x\) である。
\end{proof}

% 補題 1

\begin{lmm}\label{260422095901}
  点 \(a,b\) に対して \(a=\nu b\) となるスカラー \(\nu\) が存在しないとき、どのような点 \(x\) に対してもただ一組の数 \(\lambda,\mu\) が存在して、
  \[x=\lambda a+\mu b.\]
\end{lmm}

\begin{proof}
  順序対を用いて表せば、
  \[\langle\chi_1,\chi_2\rangle=\lambda\langle\alpha_1,\alpha_2\rangle+\mu\langle\beta_1,\beta_2\rangle.\]
  すなわち
  \[
    \left\{ \,
      \begin{aligned}
        & \alpha_1\lambda+\beta_1\mu=\chi_1, \\
        & \alpha_2\lambda+\beta_2\mu=\chi_2.
      \end{aligned}
    \right.
  \]
  ゆえに
  \[
    \lambda=\frac{\beta_2\chi_1-\beta_1\chi_2}{\alpha_1\beta_2-\alpha_2\beta_1},\quad
    \mu=\frac{\alpha_1\chi_1-\alpha_2\chi_2}{\alpha_1\beta_2-\alpha_2\beta_1}.
  \]
  ここで、\(\alpha_1\beta_2-\alpha_2\beta_1=0\) となることはあり得ない。
  もしそうだとすると、\(\langle\alpha_1,\alpha_2\rangle=\nu\langle\beta_1,\beta_2\rangle\) となる \(\nu\) が存在しないとした前提に反する。
\end{proof}

% 定理 6

\begin{thm}\label{260424194701}
  \(\overrightarrow{ab}=\nu\overrightarrow{ac}\) となるスカラー \(\nu\) が存在しないとき、どのような点 \(x\) に対してもただ一組のスカラー \(\lambda,\mu\) が存在して、
  \[\overrightarrow{ax}=\lambda\overrightarrow{ab}+\mu\overrightarrow{ac}.\]
\end{thm}

\begin{proof}
  補題 \ref{260422095901} より直ちに題意が従う。
\end{proof}


\section{角の取扱い}

% 定義 4

\begin{dfn}[\textbf{角の合同}\index{かくのごうどう@角の合同}]
  任意の点 \(x,y,z\) と \(u,v,w\) に対して、
  スカラー \(\lambda\geq 0,\ \mu\geq 0\) が存在して
  \begin{gather*}
    ||\overrightarrow{yx}||=||\lambda\overrightarrow{vu}||,\quad ||\overrightarrow{yz}||=||\mu\overrightarrow{vw}||, \\
    ||\overrightarrow{yx}-\overrightarrow{yz}||=||\lambda\overrightarrow{vu}-\mu\overrightarrow{vw}||
  \end{gather*}
  となることを、
  \(\angle xyz\equiv\angle uvw\) と表す。
\end{dfn}

% 定理 7

\begin{thm}\label{260424155001}
  \(||\overrightarrow{yx}||=||\overrightarrow{vu}||,\ ||\overrightarrow{yz}||=||\overrightarrow{vw}||\) かつ \(\angle xyz\equiv\angle uvw\) であるとき、
  任意のスカラー \(\lambda\geq 0,\ \mu\geq 0\) 対して
  \[||\lambda \overrightarrow{yx}-\mu\overrightarrow{yz}||=||\lambda \overrightarrow{vu}-\mu\overrightarrow{vw}||.\]
\end{thm}

\begin{proof}
  条件より、
  \[||\overrightarrow{yx}-\overrightarrow{yz}||^2=||\overrightarrow{vu}-\overrightarrow{vw}||^2.\]
  すなわち
  \[\overrightarrow{yx}\cdot\overrightarrow{yz}=\overrightarrow{vu}\cdot\overrightarrow{vw}.\]

  ゆえに
  \begin{gather*}
    ||\lambda\overrightarrow{yx}-\mu\overrightarrow{yz}||^2-||\lambda\overrightarrow{vu}-\mu\overrightarrow{vw}||^2 \\
    =2\lambda\mu(\overrightarrow{yx}\cdot\overrightarrow{yz})-2\lambda\mu(\overrightarrow{vu}\cdot\overrightarrow{vw})=0.
  \end{gather*}

  よって、定理の結論が示された。
\end{proof}

\begin{col*}
  定理 \ref{260424155001} により、角の合同が同一律・反射律だけでなく推移律も満たすことがわかる。
\end{col*}

% 定義 5

\begin{dfn}[\textbf{角の大小}\index{かくのだいしょう@角の大小}]
  任意の点 \(x,y,z\) と \(u,v,w\) に対して、
  \[\angle xyz\equiv\angle x'y'z',\quad \angle uvw\equiv\angle u'v'w'\]
  となる点 \(x',y',z'\) と \(u',v',w'\) が存在し（同じでもよい）、
    スカラー \(\lambda\geq 0,\ \nu>1\) が存在して
  \begin{gather*}
    ||\overrightarrow{y'z'}||=||\overrightarrow{v'w'}||, \\
    \overrightarrow{y'x'}-\overrightarrow{y'z'}=\nu(\lambda\overrightarrow{v'u'}-\overrightarrow{v'w'})
  \end{gather*}
  となることを、
  \[\angle xyz>\angle uvw\]
  と表す。
\end{dfn}

\begin{col*}
  \(\lambda=0\) とすれば、\(180\) 度の角が最大だとわかる。
  （それより大きい角を扱うには、別種の角として定義する必要がある。）
\end{col*}

  \chapter{定理}

% 定義 6

\begin{dfn}[\textbf{直線上}\index{ちょくせんじょう@直線上}]
  点 \(x\) は直線 \(ab\) 上にある、を
  \[\exists\lambda\{x-a=\lambda(b-a)\}\]
  と定義する。

  点 \(x\) が直線 \(yz\) 上にあることを、\(y,x,z\) は\textbf{一直線上}\index{いっちょくせんじょう@一直線上}にある、ともいう。また、このことを、直線 \(yz\) は点 \(x\) を\textbf{通る}\index{とおる@通る}、ともいう。
\end{dfn}

% 定義 7

\begin{dfn}[\textbf{半直線上}\index{はんちょくせんじょう@半直線上}]
  点 \(x\) は半直線 \(ab\) 上にある、を
  \[\exists\lambda\{x-a=\lambda(b-a)\land\lambda\geq 0\}\]
  と定義する。
\end{dfn}

% 定義 8

\begin{dfn}[\textbf{半直線の延長上}\index{はんちょくせんのえんちょうじょう@半直線の延長上}]
  点 \(x\) は半直線 \(ab\) の延長上にある、を
  \[\exists\lambda\{x-a=\lambda(b-a)\land\lambda<0\}\]
  と定義する。
\end{dfn}

% 定義 9

\begin{dfn}[\textbf{線分上}\index{せんぶんじょう@線分上}]
  点 \(x\) は線分 \(ab\) 上にある、を
  \[\exists\lambda\{x-a=\lambda(b-a)\land0\leq\lambda\leq 1\}\]
  と定義する。
\end{dfn}

% 定義 10

\begin{dfn}[\textbf{交わる}\index{まじわる@交わる}・\textbf{交点}\index{こうてん@交点}]
  点 \(x\) が直線 \(ab\) 上にあり、直線 \(cd\) 上にもあるとき、直線 \(ab\) と直線 \(cd\) は交わるといい、\(x\) を直線 \(ab\) と直線 \(cd\) の交点という。

  直線、半直線、線分と直線、半直線、線分の組み合わせにおける他の場合も同様である。
\end{dfn}

% 定義 11

\begin{dfn}[\textbf{直線図形の合同}\index{ちょくせんずけいのごうどう@直線図形の合同}]
  任意の点について、それが直線 \(ab\) 上にあることと、直線 \(cd\) 上にあることが同値なとき、直線 \(ab\) と直線 \(cd\) は合同であるという。

  半直線や線分の場合も同様である。

  この合同を等しいなどとも言い表す。
\end{dfn}

\begin{col*}
  直線・半直線・線分の合同は、同値関係の性質を引き継ぐので、明らかに同一律・反射律・推移律が成り立つ。

  また、「点 \(z\) は直線 \(xy\) 上にある」かつ「直線 \(xy\) と直線 \(uv\) は合同」ならば、「点 \(z\) は直線 \(uv\) 上にある」。
\end{col*}

% 定義 12

\begin{dfn}[\textbf{関数 \(d\)}\index{かんすうでー@関数 \(d\)}・\textbf{長さ}\index{ながさ@長さ}]
  ベクトルとは限らない任意の \(x,y\) に対して、関数 \(d\) を
  \[d(x,y)=||y-x||\]
  と定義する。

  \(x,y\) が点であるとき、\(d(x,y)\) を線分 \(xy\) の長さともいい、
  \[\overline{xy}\]
  と表す。
\end{dfn}

% 定義 13

\begin{dfn}[\textbf{半平面上}\index{はんへいめん@半平面}・\textbf{へり}\index{へり@へり}・\textbf{同じ側}\index{おなじがわ@同じ側}]
  点 \(a,b,c\) が一直線上にないとき、点 \(x\) は半平面 \(cab\) の上にある、を
  \[\exists\lambda\mu\{x-a=\lambda(b-a)+\mu(c-a)\land\mu\geq 0\}\]
  と定義する。
  点 \(x\) は直線 \(ab\) をへりとし点 \(c\) を含む半平面上にある、ともいう。

  また、このとき \(\mu\neq 0\) ならば、点 \(x\) と点 \(c\) は直線 \(ab\) に関して同じ側にある、という。
\end{dfn}

% 定義 14

\begin{dfn}[\textbf{共役な半平面上}\index{きょうやくなはんへいめん@共役な半平面上}・\textbf{へり}・\textbf{反対側}\index{はんたいがわ@反対側}]\label{260502235601}
  点 \(a,b,c\) が一直線上にないとき、点 \(x\) は共役な半平面 \(cab\) の上にある、を
  \[\exists\lambda\mu\{x-a=\lambda(b-a)+\mu(c-a)\land\mu\leq 0\}\]
  と定義する。
  点 \(x\) は直線 \(ab\) をへりとし点 \(c\) を含む共役な半平面上にある、ともいう。

  また、このとき \(\mu\neq 0\) ならば、点 \(x\) と点 \(c\) は直線 \(ab\) に関して反対側にある、という。
\end{dfn}

% 定理 8

\begin{thm}\label{360424013601}
  与えられた点と同じ側にある点とを結ぶ線分は与えられた直線と両端以外で交わらない。
\end{thm}

\begin{proof}
  点 \(p\) が半平面 \(cab\) 上にあって、直線 \(ab\) 上にないとする。

  半平面の定義より
  \[p-a=\lambda(b-a)+\mu(c-a),\quad \mu>0.\]

  さて、線分 \(cp\) 上の両端以外の任意の点 \(x\) について、
  \[x-c=\nu(p-c),\quad 0<\nu<1.\]
  また
  \[a=\nu a+(1-\nu)a\]
  であるから、
  \begin{align*}
    x-a &=\nu(p-a)+(1-\nu)(c-a) \\
        &=\nu\{\lambda(b-a)+\mu(c-a)\}+(1-\nu)(c-a) \\
        &=\lambda\nu(b-a)+(\mu\nu+1-\nu)(c-a).
  \end{align*}

  \(\mu\nu+1-\nu>0\) なので、点 \(x\) は直線 \(ab\) 上にない。
\end{proof}

% 定理 9

\begin{thm}\label{260424012501}
  与えられた点と反対側にある点とを結ぶ線分は与えられた直線と両端以外で交わる。
\end{thm}

\begin{proof}
  点 \(p\) が半平面 \(cab\) 上にあって、直線 \(ab\) 上にないとする。

  半平面の定義より
  \[p-a=\lambda(b-a)+\mu(c-a),\quad \mu<0.\]

  さて、線分 \(cp\) 上の両端以外の任意の点 \(x\) について、
  \[x-c=\nu(p-c),\quad 0<\nu<1.\]
  また
  \[a=\nu a+(1-\nu)a\]
  であるから、
  \begin{align*}
    x-a &=\nu(p-a)+(1-\nu)(c-a) \\
        &=\nu\{\lambda(b-a)+\mu(c-a)\}+(1-\nu)(c-a) \\
        &=\lambda\nu(b-a)+(\mu\nu+1-\nu)(c-a).
  \end{align*}

  \(\mu\nu+1-\nu=0\) ならば、点 \(x\) は直線 \(ab\) 上にある。
  \[\nu=\frac{-1}{\mu-1}\]
  であるが、\(\mu<0\) より \(0<\nu<1\) となるので、正当である。
\end{proof}

% 補題 2

\begin{lmm}\label{260502235101}
  点 \(d\) が半平面 \(cab\) 上にないとき、半平面 \(cab\) は共役な半平面 \(dab\) と合同（定義は直線図形のそれと同様）である。
\end{lmm}

\begin{proof}
  「\(x\) は共役な半平面 \(dab\) 上にある」が「\(x\) は半平面 \(cab\) 上にある」と同値であることを言えばよい。

  \(d\) の満たすべき条件は
  \[d-a=\alpha(b-a)+\beta(c-a),\quad \beta<0\]
  である。
  ゆえに、\(x\) が共役な半平面 \(dab\) 上にあるということは、\(\mu\leq 0\) として
  \begin{align*}
    x-a &=\lambda(b-a)+\mu(d-a) \\
        &=\lambda(b-a)+\mu\{\alpha(b-a)+\beta(c-a)\} \\
        &=(\lambda+\alpha\mu)(b-a)+\beta\mu(c-a)
  \end{align*}
  と表される。
  \(\lambda+\alpha\mu\) の値は \(\mu\) を固定したとしてもあらゆるスカラーを渡る。
  \(\beta\mu\) の値は \(\beta<0\) だからあらゆる非負のスカラーを渡る。

  ゆえに、この等式は \(\lambda',\mu'\) を任意のスカラーとして、
  \[x-a=\lambda'(b-a)+\mu'(c-a),\quad \mu'\geq 0\]
  と表されるが、これは \(x\) が半平面 \(cab\) 上にあることの定義にほかならない。
\end{proof}

% 定理 10

\begin{thm}[\textbf{Pasch の公理}\index{ぱっしゅのこうり@Pasch の公理}]
  三角形の一辺と交わって両端の頂点を通らない直線は、その辺に対する頂点を通るか、残りのどちらか一方の辺と交わる。
\end{thm}

\begin{proof}
  三角形 \(abc\) において直線 \(de\) が辺 \(bc\) と点 \(d\) で交わり、\(d\neq b,\ d\neq c\) とする。

  定理 \ref{360424013601} の対偶より、\(b\) と \(c\) は直線 \(de\) の互いに反対側にある。
  そして、定理 \ref{260424194701} と、半平面および共役な半平面との定義より、頂点 \(a\) は、直線 \(de\) 上にあるか、さもなければ \(de\) をへりとする二つの半平面のどちらか一方の上にのみある。
  すなわち、半平面 \(bde\) か共役な半平面 \(bde\) である。

  ところが \(c\) は直線 \(de\) に関して \(b\) の反対側にあるのだから、「反対側にある」の定義（\ref{260502235601}）より、共役な半平面 \(bde\) 上にあるのであり、かつ \(c\) は直線 \(de\) 上にないので、これは \(c\) が半平面 \(bde\) 上にないことを意味する。
  すなわち、補題 \ref{260502235101} より半平面 \(bde\) は共役な半平面 \(cde\) と合同である。

  \(a\) が共役な半平面 \(bde\) 上の点なら、\(a\) と \(b\) は互いに反対側にある。
  ゆえに定理 \ref{260424012501}より、直線 \(de\) は線分 \(ab\) と交わる。

  \(a\) が共役な半平面 \(cde\) 上の点なら、\(a\) と \(c\) は互いに反対側にある。
  ゆえに定理 \ref{260424012501}より、直線 \(de\) は線分 \(ac\) と交わる。
\end{proof}

% 定義 15

\begin{dfn}[\textbf{正射影}\index{せいしゃえい@正射影}]
  任意の点 \(x\neq\langle0,0\rangle\) に対して、
  \[\frac{x\cdot y}{||x||^2}x\]
  を \(x\) の上への \(y\) の正射影という。
\end{dfn}

% 定理 11

\begin{thm}\label{260427130101}
  \(x\) の上への \(y\) の正射影を \(y_0\) とするとき、
  \[x\cdot(y-y_0)=0\]
\end{thm}

\begin{proof}
  \begin{align*}
    x\cdot(y-y_0) &=x\cdot(y-\frac{x\cdot y}{||x||^2}x)=x\cdot y-x\cdot\frac{x\cdot y}{||x||^2}x \\
                  &=x\cdot y-\frac{x\cdot y}{||x||^2}(x\cdot x)=x\cdot y-\frac{x\cdot y}{||x||^2}||x||^2 \\
                  &=x\cdot y-x\cdot y=0.
  \end{align*}
\end{proof}

% 定義 16

\begin{dfn}[\textbf{ベクトルの垂直}\index{すいちょく@垂直（ベクトルの）}]
  ベクトル \(x,y\) について \(x\cdot y=0\) であることを、\(x\) と \(y\) は垂直である、\(x\) は \(y\) に垂直である、などという。
\end{dfn}

% 定義 17

\begin{dfn}[\textbf{ベクトルの平行}\index{へいこう@平行（ベクトルの）}]
  ベクトル \(x,y\neq\langle0,0\rangle\) に対し、スカラー \(\lambda\) が存在して \(x=\lambda y\) となることを、\(x\) と \(y\) は平行である、\(x\) は \(y\) に平行である、などという。
\end{dfn}

% 定義 18

\begin{dfn}[\textbf{三角形の合同}\index{さんかくけいのごうどう@三角形の合同}]
  任意の三角形 \(xyz\) と \(uvw\) について、
  \begin{align*}
    & \overline{yz}=\overline{vw}, && \overline{zx}=\overline{wu}, && \overline{xy}=\overline{uv}, \\
    & \angle zxy\equiv\angle wuv, && \angle xyz\equiv\angle uvw, && \angle yzx\equiv\angle vwu
  \end{align*}
  であることを、三角形 \(xyz\) と三角形 \(uvw\) は合同である、といい、
  \[\triangle xyz\equiv\triangle uvw\]
  と書く。
\end{dfn}

\begin{col*}
  この体系においては、三角形の角の合同は頂点に対する辺（すなわち線分）の合同を意味するので、この定義は実質的に辺の合同のみで成立する。
\end{col*}

% 補題 3

\begin{lmm}\label{260424041701}
  \[x\cdot y=0\rightarrow||x-y||^2=||x||^2+||y||^2.\]
\end{lmm}

\begin{proof}
  \(||x-y||^2\) を展開すれば明らか。
\end{proof}

% 補題 4

\begin{lmm}\label{260424041702}
  半直線 \(ab\) と線分 \(cd\) が与えられたとき、半直線 \(ab\) 上に点 \(x\) が存在して
  \[\overline{ax}=\overline{cd}.\]
\end{lmm}

\begin{proof}
  半直線 \(ab\) 上の点 \(x\) を表す式 \(x-a=\lambda(b-a)\) において \(\displaystyle \lambda=\frac{\overline{cd}}{||b-a||}\) とおくと、
  \[||x-a||=\frac{\overline{cd}}{||b-a||}||b-a||=\overline{cd}.\]
\end{proof}

% 補題 5

\begin{lmm}\label{260424041703}
  どのような直線とその上にあるどの点に対しても、その点から出る垂直なベクトルがある。
\end{lmm}

\begin{proof}
  まず、どのようなベクトルにもそれに垂直なベクトルがあることを示す。

  \(x=\langle\alpha,\beta\rangle\) であるとすると、\(y=\langle\beta,-\alpha\rangle\) とおけば明らかに \(x\cdot y=x\cdot(-y)=0\) である。

  ゆえに、直線のベクトルに垂直なベクトル \(n\) を選び、与えられた点 \(a\) に対して、\(b=n+a\) としてベクトル \(b\) を求めれば、それが求めるべきものである。
\end{proof}

% 準公理 1

\begin{semi}[\textbf{三角形の移動の公理}]
  半直線とその上にも延長上にもない点および三角形が与えられたとき、
  半直線の上に点をとり、さらに、半直線とその延長をへりとする半平面のうち、与えられた点を含む側の上に点をとり、与えられた三角形と合同な三角形を作ることができる。
\end{semi}

\begin{proof}
  与えられた三点を \(a,b,c\) とし、三角形 \(xyz\) が与えられているとする。

  \(\overrightarrow{yz}\) の上への \(\overrightarrow{yx}\) の正射影を \(\overrightarrow{yx}_0\) とすると、定理 \ref{260427130101} より
  \[\overrightarrow{yz}\cdot\overrightarrow{x_0x}=\overrightarrow{yz}\cdot(\overrightarrow{yx}-\overrightarrow{yx}_0)=0.\]

  補題 \ref{260424041701}より
  \[||\overrightarrow{x_0y}||^2+||\overrightarrow{x_0x}||^2=||\overrightarrow{x_0y}-\overrightarrow{x_0x}||^2=||\overrightarrow{xy}||^2.\]

  同様に
  \[||\overrightarrow{x_0z}||^2+||\overrightarrow{x_0x}||^2=||\overrightarrow{x_0z}-\overrightarrow{x_0x}||^2=||\overrightarrow{xz}||^2.\]

  さて、補題 \ref{260424041702} により、半直線 \(bc\) 上に \(ba_0=yx_0\) となる点 \(a_0\) と、\(bc'=yz\) となる点 \(c'\) をとり、補題 \ref{260424041703} により、\(a_0\) から出る垂直なベクトルを立てる。
  そして、そのベクトルによって定まる直線上の点として \(a_0a'=x_0x\) となる点 \(a'\) を、ただし直線 \(bc\) に関して点 \(a\) と同じ側にあるようにとれば、
  \[||\overrightarrow{a_0b}||^2+||\overrightarrow{a_0a'}||^2=||\overrightarrow{a_0b}-\overrightarrow{a_0a'}||^2=||\overrightarrow{a'b}||^2.\]
  \[||\overrightarrow{a_0c'}||^2+||\overrightarrow{a_0a'}||^2=||\overrightarrow{a_0c'}-\overrightarrow{a_0a'}||^2=||\overrightarrow{a'c'}||^2.\]
  ゆえに
  \[a'b=xy,\quad a'c'=xz.\]
  \(bc'=yz\) は既出である。
\end{proof}

% 補題 6

\begin{lmm}\label{260424111401}
  異なる二点 \(x,y\) が共に直線 \(ab\) 上にあるならば、直線 \(xy\) と直線 \(ab\) は合同である。
\end{lmm}

\begin{proof}
  条件より
  \[x-a=\lambda_1(b-a),\quad y-a=\lambda_2(b-a).\]
  ゆえに
  \[y-x=(\lambda_2-\lambda_1)(b-a).\]

  ここで、
  \[z-x=\lambda(y-x)\]
  とするならば、
  \[z-x=\lambda(\lambda_2-\lambda_1)(b-a)\]
  となるので、最初の等式より
  \[z-a=\{\lambda(\lambda_2-\lambda_1)+\lambda_1\}(b-a).\]

  逆に、
  \[z-a=\lambda'(b-a)\]
  とする。
  \(x,y\) が異なる二点であることより \(\lambda_1\neq\lambda_2\) である。
  ゆえに、
  \[\lambda'=\lambda(\lambda_2-\lambda_1)+\lambda_1\]
  となる \(\lambda\) は必ず存在するので、先ほどの推論を逆向きにたどれば
  \[z-x=\lambda(y-x).\]
\end{proof}

% 準公理 2

\begin{semi}[\textbf{直線の唯一存在の公理}]
  異なる二点を通る直線はただ一つある。
\end{semi}

\begin{proof}
  まず、高々一つであることを示す。

  異なる二点 \(a,b\) を通る直線は、それがどのような直線、例えば直線 \(xy\) であれ、\(a,b\) が直線 \(xy\) 上にあれば補題 \ref{260424111401} により直線 \(ab\) と直線 \(xy\) は合同である。

  すなわち、直線の合同が推移律を満たすことより、異なる二点 \(a,b\) を通る直線はすべて合同である。

  直線が存在することに関しては、この体系において直線は対象物でないので問題を正面から扱うことができないが、「その直線上の点が存在する」と解すれば、直線は二点で規定され、少なくともその二点を通るのだから、自明である。
\end{proof}

% 定理 12

\begin{thm}[\textbf{Schwarz の不等式}\index{しゅわるつのふとうしき@Schwarz の不等式}]\label{260424234101}
  \[|x\cdot y|\leq||x||||y||.\]
\end{thm}

\begin{proof}
  \(x=\langle0,0\rangle\) ならば明らかなので、\(x\neq\langle0,0\rangle\) とする。
  内積の性質より、
  \[(\lambda x+y)\cdot(\lambda x+y)\geq 0.\]
  左辺を展開すると、
  \[(x\cdot x)\lambda^2+2(x\cdot y)\lambda+(y\cdot y)\geq 0.\]
  \(x\neq\langle0,0\rangle\) より \(x\cdot x>0\) であるから、左辺を平方完成すると
  \[(x\cdot x)\left(\lambda+\frac{x\cdot y}{x\cdot x}\right)^2-\frac{(x\cdot y)^2-(x\cdot x)(y\cdot y)}{x\cdot x}\geq 0.\]
  ゆえに、左辺は \(\displaystyle \lambda=-\frac{x\cdot y}{x\cdot x}\) のとき最小値をとり、それが \(\geq 0\) でなければならない。
  よって
  \[(x\cdot y)^2\leq(x\cdot x)(y\cdot y).\]
  ゆえに
  \[|x\cdot x|\leq||x||||y||.\]
\end{proof}

% 補題 7

\begin{lmm}\label{260424225402}
  \[||x+y||\leq||x||+||y||.\]
\end{lmm}

\begin{proof}
  左辺を展開して定理 \ref{260424234101} を用いれば示せる。
\end{proof}

% 補題 8

\begin{lmm}\label{260424225401}
  \(x\neq y,\ z\neq y\) である任意の三点について、
  \[\overline{yx}+\overline{yz}=\overline{xz}\rightarrow\exists\lambda(\overrightarrow{yz}=-\lambda\overrightarrow{yx}\land\lambda>0).\]
\end{lmm}

\begin{proof}
  \(\overline{yx}+\overline{yz}=\overline{xz}\) より、
  \[||\overrightarrow{yx}||+||\overrightarrow{yz}||=||\overrightarrow{yz}-\overrightarrow{yx}||.\]

  両辺を自乗して整理すると、
  \(||\overrightarrow{yx}||||\overrightarrow{yz}||=-\overrightarrow{yx}\cdot\overrightarrow{yz}.\)

  さて、\(||\overrightarrow{yz}+\lambda\overrightarrow{yx}||=0\) とおくと、

  \begin{align*}
    0 &=||\overrightarrow{yz}+\lambda\overrightarrow{yx}||^2=(\overrightarrow{yz}+\lambda\overrightarrow{yx})\cdot(\overrightarrow{yz}+\lambda\overrightarrow{yx}) \\
      &=(\overrightarrow{yx}\cdot\overrightarrow{yx})\lambda^2+2(\overrightarrow{yx}\cdot\overrightarrow{yz})\lambda+\overrightarrow{yz}\cdot\overrightarrow{yz} \\
      &=||\overrightarrow{yx}||^2\lambda^2-2||\overrightarrow{yx}||||\overrightarrow{yz}||\lambda+||\overrightarrow{yz}||^2 \\
      &=(||\overrightarrow{yx}||\lambda-||\overrightarrow{yz}||)^2.
  \end{align*}

  ゆえに、\(\displaystyle \lambda=\frac{||\overrightarrow{yz}||}{||\overrightarrow{yx}||}\,(>0)\) について \(\overrightarrow{yz}+\lambda\overrightarrow{yx}=\langle0,0\rangle.\)
\end{proof}

% 準公理 3

\begin{semi}[\textbf{線分の最短性の公理}]
  三角形の二辺の長さの和は残り一辺の長さよりも大きい。
\end{semi}

\begin{proof}
  三角形の頂点を \(a,b,c\) とすると、補題 \ref{260424225402} より
  \[||\overrightarrow{ac}||=||\overrightarrow{ab}+\overrightarrow{bc}||\leq||\overrightarrow{ab}||+||\overrightarrow{bc}||.\]

  補題 \ref{260424225401} より、\(\overline{ab}+\overline{bc}=\overline{ac}\) は成り立たない。
\end{proof}

% 補題 9

\begin{lmm}\label{260424203201}
  異なる二つの直線のベクトルが平行でないならば、それらは交わる。
\end{lmm}

\begin{proof}
  ベクトル \(\overrightarrow{ab},\overrightarrow{cd}\) が平行でないとし、\(\overrightarrow{ae}=\overrightarrow{cd}\) となる点 \(e\) をとる。

  \(\overrightarrow{ab}\) と \(\overrightarrow{ac}\) が平行であるとすると、点 \(c\) は直線 \(ab\) 上の点であり、二つの直線は \(c\) で交わることになって、証明は終わりである。
  ゆえに、平行ではないとする。

  このとき、定理 \ref{260424194701} より \(\overrightarrow{ae}=\lambda\overrightarrow{ab}+\mu\overrightarrow{ac}\) となるスカラー \(\lambda,\mu\) がある。

  さて、ここで \(\mu=0\) とすると、\(\overrightarrow{cd}=\overrightarrow{ae}=\lambda\overrightarrow{ab}\) より、\(\overrightarrow{cd}\) と \(\overrightarrow{ab}\) が平行になってしまうので、\(\mu\neq 0.\)

  ゆえに、\(\displaystyle \overrightarrow{af}=-\frac{\lambda}{\mu}\overrightarrow{ab}\) とおくことができ、
  \begin{align*}
    \overrightarrow{cf} &=\overrightarrow{af}-\overrightarrow{ac}=-\frac{\lambda}{\mu}\overrightarrow{ab}-\overrightarrow{ac} \\
                        &=-\frac{1}{\mu}(\lambda\overrightarrow{ab}+\mu\overrightarrow{ac}) \\
                        &=-\frac{1}{\mu}\overrightarrow{ae}=-\frac{1}{\mu}\overrightarrow{cd}.
  \end{align*}
  これは、直線 \(ab\) 上の点として定義した点 \(f\) が直線 \(cd\) 上にもあることを示している。
\end{proof}

% 準公理 4

\begin{semi}[\textbf{平行線の唯一性の公理}]
  直線外の一点を通ってその直線に決して交わらない直線は高々一つ。
\end{semi}

\begin{proof}
  その直線外の一点が合同でない二つの直線上にあったとしよう。
  補題 \ref{260424203201} より、それら二つの直線のそれぞれのベクトルは与えられた直線のベクトルと平行である。
  与えられた点を \(a\) とすれば、
  \[\overrightarrow{ab}=\lambda_1\overrightarrow{xy},\quad \overrightarrow{ac}=\lambda_2\overrightarrow{xy}.\]

  ゆえに
  \[\lambda_2\overrightarrow{ab}=\lambda_1\overrightarrow{ac}.\]
  ところが、この等式を用いると一方の直線上の点が必ず他方の直線上にもあることが分かるので、合同でない二つの直線という前提に反する。
\end{proof}

これで準公理が出揃ったので、図を見れば明らかなことについてはいちいち証明せず、準公理のみを用いる方針で行くことができる。

  \input{appendix.tex}
  \begin{thebibliography}{99}
    \bibitem{seki} 赤攝也.『現代の初等幾何学』(ちくま学芸文庫). 筑摩書房, 2019.
    \bibitem{kodaira} 小平邦彦.『幾何への誘い』(岩波現代文庫). 岩波, 2010.
  \end{thebibliography}
  \clearpage
  \printindex
\end{document}
